\section{Animação, FPS e Interatividade}

\begin{frame}{Como Criar a Ilusão de Movimento?}
    \begin{columns}
        \begin{column}{0.5\textwidth}
            \textbf{A Lógica da Física:}
            \begin{itemize}
                \item Uma animação consiste em mudar levemente as coordenadas $(x, y)$ de um objeto a cada passagem pelo laço principal.
                \item Equação básica da velocidade constante:
                \[
                    x \leftarrow x + v_x
                \]
                \[
                    y \leftarrow y + v_y
                \]
            \end{itemize}
        \end{column}
        \begin{column}{0.5\textwidth}
            \begin{alertblock}{O Erro Clássico do Rastro}
                Se você esquecer de chamar \texttt{tela.fill(COR)} antes de desenhar o novo quadro, a tela não será limpa!
                \vspace{0.2cm}
                
                O objeto desenhado no quadro anterior continuará lá, criando um rastro borrado infinito.
            \end{alertblock}
        \end{column}
    \end{columns}
\end{frame}

\begin{frame}{O Papel Crucial de \texttt{pygame.time.Clock()}}
    \begin{block}{O que acontece sem o controle de FPS?}
        \begin{itemize}
            \item O computador rodará o \texttt{while} na velocidade máxima que o processador aguentar (milhares de ciclos por segundo).
            \item O jogo consumirá \textbf{100\% de CPU}, e rodará ultra rápido em máquinas potentes e lento em computadores antigos!
        \end{itemize}
    \end{block}

    \vspace{0.3cm}
    \begin{columns}
        \begin{column}{0.5\textwidth}
            \begin{block}{Com \texttt{relogio.tick(60)}}
                \begin{itemize}
                    \item O Pygame calcula quanto tempo o quadro demorou e coloca a CPU para dormir (\textit{sleep}) pelo tempo restante.
                    \item Garante que o jogo rode a exatos \textbf{60 FPS} em qualquer computador.
                \end{itemize}
            \end{block}
        \end{column}
        \begin{column}{0.5\textwidth}
            \textbf{Fórmula do Tempo de Quadro:}
            \[
                \Delta t = \frac{1000\text{ ms}}{60\text{ FPS}} \approx 16{,}6\text{ ms por quadro}
            \]
        \end{column}
    \end{columns}
\end{frame}

\begin{frame}[fragile]{Movimentação Interativa pelo Teclado}
    Para movimentar personagens de forma contínua e suave, utilizamos \texttt{pygame.key.get\_pressed()}:

    \begin{lstlisting}[language=Python]
# Captura o estado continuo de todas as teclas
teclas = pygame.key.get_pressed()

# Movimento Horizontal (Setas ou A / D)
if teclas[pygame.K_LEFT] or teclas[pygame.K_a]:
    jogador_x -= velocidade
if teclas[pygame.K_RIGHT] or teclas[pygame.K_d]:
    jogador_x += velocidade

# Movimento Vertical (Setas ou W / S)
if teclas[pygame.K_UP] or teclas[pygame.K_w]:
    jogador_y -= velocidade
if teclas[pygame.K_DOWN] or teclas[pygame.K_s]:
    jogador_y += velocidade
    \end{lstlisting}

    \vspace{0.2cm}
    \textbf{Diferença Fundamental:}
    \begin{itemize}
        \item \texttt{pygame.KEYDOWN} (na fila de eventos): Dispara \textbf{uma única vez} quando a tecla desce (ótimo para tiro ou pulo).
        \item \texttt{pygame.key.get\_pressed()}: Responde \textbf{enquanto a tecla estiver pressionada} (ótimo para corrida).
    \end{itemize}
\end{frame}
